谱曲面的两处无穷远分别为 $z=\infty$ 与 $z=0$. 以下在固定 $t$ 下推导修正 Jost 解及透射系数的展开. 为使一阶系数及二阶余项均可控制, 假设
\begin{equation}\label{eq:highenergy-assumption}
q-q_0\in W^{2,1}(\mathbb R).
\end{equation}
此时 $q$ 有界, $q_x$ 连续且在两端趋于零, $q_x,q_{xx}\in L^1(\mathbb R)$.

\begin{proposition}[两处无穷远的渐近展开]\label{prop:highenergy}
在条件 \eqref{eq:highenergy-assumption} 下, 记
\begin{equation}\label{eq:asymptotic-coefficients}
B(q)=\sg Q=\begin{pmatrix}0&q\\\conj{q}&0\end{pmatrix},\qquad
J_\pm(x)=\int_{\pm\infty}^{x}(q_0^2-|q(y)|^2)\,\dd y.
\end{equation}
各 Jost 解析列在相应闭域内满足
\begin{equation}\label{eq:mu-highenergy}
\mu_\pm(x,z)=I+\dfrac{A_\pm(x)}{z}+O(z^{-2}),\qquad
A_\pm=-\ii B(q)+\ii J_\pm\sg,\qquad z\to\infty,
\end{equation}
余项对实 $x$ 一致, 矩阵式按各自解析列理解. 令
\[
J_{\rm tot}=\int_{-\infty}^{+\infty}(q_0^2-|q|^2)\,\dd x,
\]
则透射系数满足
\begin{equation}\label{eq:transmission-highenergy}
\begin{aligned}
s_{22}(z)&=1+\dfrac{\ii J_{\rm tot}}{z}+O(z^{-2}),&&z\to\infty,\quad z\in\overline{D_+},\\
s_{11}(z)&=1-\dfrac{\ii J_{\rm tot}}{z}+O(z^{-2}),&&z\to\infty,\quad z\in\overline{D_-}.
\end{aligned}
\end{equation}
在另一处无穷远,
\begin{equation}\label{eq:smallmu}
\mu_\pm(x,z)=-\dfrac{\ii q_0}{z}\sigma_1
+\dfrac{1}{q_0}\begin{pmatrix}q(x)&-J_\pm(x)\\J_\pm(x)&\conj{q(x)}\end{pmatrix}
+O(z),\qquad z\to0,
\end{equation}
以及
\begin{equation}\label{eq:transmission-smallenergy}
\begin{aligned}
s_{22}(z)&=1+\dfrac{\ii J_{\rm tot}}{q_0^2}z+O(z^2),&&z\to0,\quad z\in\overline{D_+},\\
s_{11}(z)&=1-\dfrac{\ii J_{\rm tot}}{q_0^2}z+O(z^2),&&z\to0,\quad z\in\overline{D_-}.
\end{aligned}
\end{equation}
完整散射矩阵在实轴两端满足
\begin{equation}\label{eq:scattering-two-infinities}
\begin{aligned}
S(z)&=I-\dfrac{\ii J_{\rm tot}}{z}\sg+O(z^{-2}),&&z\to\infty,\\
S(z)&=I-\dfrac{\ii J_{\rm tot}}{q_0^2}z\sg+O(z^2),&&z\to0.
\end{aligned}
\end{equation}
\end{proposition}
\begin{proof}
先确定展开的系数. 将空间方程写为
\begin{equation}\label{eq:mu-highenergy-system}
\mu_x=-\dfrac{\ii z}{2}[\sg,\mu]+Q\mu
+\dfrac{\ii q_0^2}{2z}(\sg\mu+\mu\sg).
\end{equation}
以 $\mu=I+\dfrac{A}{z}+\dfrac{A_2}{z^2}+\cdots$ 为待证展开, 比较常数项与 $z^{-1}$ 项, 分别得到
\[
-\dfrac{\ii}{2}[\sg,A]+Q=0,\qquad
A_x=-\dfrac{\ii}{2}[\sg,A_2]+QA+\ii q_0^2\sg.
\]
第一式确定 $A$ 的非对角部分为 $-\ii\sg Q=-\ii B(q)$. 对第二式取对角部分, 使用 $QB=|q|^2\sg$, 得到
\[
(A_{\rm d})_x=\ii(q_0^2-|q|^2)\sg,\qquad
A_{\rm d}(\pm\infty)=0,
\]
故 $A_\pm=-\ii B+\ii J_\pm\sg$. 以下由积分方程证明此展开及其一致余项, 无须预先假定渐近级数可逐项求导.

作局部矩阵变换
\[
\mathcal T(x,z)=I-\dfrac{\ii B(q(x))}{z},\qquad \mu_\pm=\mathcal T n_\pm.
\]
由 $B\sg=-\sg B$, $B^2=|q|^2I$,
\[
\det\mathcal T=1+\dfrac{|q|^2}{z^2},\qquad
\mathcal T^{-1}=\dfrac{I+\dfrac{\ii B}{z}}{1+\dfrac{|q|^2}{z^2}}.
\]
当 $|z|$ 足够大时, $\mathcal T,\mathcal T^{-1}$ 对所有实 $x$ 一致有界; 在归一化端 $\mathcal T\to Y$, 故 $n_\pm\to I$. 利用 $k+\lambda=z$ 和 $\lambda-k=\dfrac{q_0^2}{z}$, 有
\begin{align*}
U\mathcal T+\ii\lambda\mathcal T\sg-\mathcal T_x
&=\ii(\lambda-k)\sg+Q-\dfrac{k+\lambda}{z}\sg B
-\dfrac{\ii}{z}QB+\dfrac{\ii}{z}B_x\\
&=\dfrac{\ii}{z}\bigl((q_0^2-|q|^2)\sg+B_x\bigr).
\end{align*}
因此
\begin{equation}\label{eq:n-volterra-system}
(n_\pm)_x=-\ii\lambda[\sg,n_\pm]+\mathcal Cn_\pm,\qquad
\mathcal C=\dfrac{1}{z}\mathcal T^{-1}\mathcal C_1,\qquad
\mathcal C_1=\ii\bigl((q_0^2-|q|^2)\sg+B_x\bigr).
\end{equation}
由 \eqref{eq:highenergy-assumption}, $q_0^2-|q|^2,q_x\in L^1$, 从而
\begin{equation}\label{eq:asymptotic-L1-control}
\|\mathcal C(\cdot,z)\|_{L^1}\le\dfrac{C}{|z|},\qquad
\left\|\mathcal C(\cdot,z)-\dfrac{\mathcal C_1}{z}\right\|_{L^1}
\le\dfrac{C}{|z|^2}.
\end{equation}
第 $j$ 列的 Volterra 方程为
\[
n_{\pm,j}(x,z)=e_j+\int_{\pm\infty}^{x}
e^{-\ii\lambda(x-y)(\sg-s_jI)}\mathcal C(y,z)n_{\pm,j}(y,z)\,\dd y.
\]
其中 $s_1=1,s_2=-1$, 传播核在相应闭解析域及积分方向上范数不超过 1. 与定理 \ref{thm:jost} 相同的有序积分估计给出第 $l$ 次迭代项的界 $\dfrac{C^l}{l!|z|^l}$, 所以
\[
\|n_\pm-I\|_{L^\infty_x}\le e^{\frac{C}{|z|}}-1=O(z^{-1}),\qquad
\sum\limits_{l=2}^{\infty}\dfrac{C^l}{l!|z|^l}=O(z^{-2}).
\]
结合 \eqref{eq:asymptotic-L1-control}, 只需计算以 $\dfrac{\mathcal C_1}{z}$ 代替 $\mathcal C$ 的第一次迭代. 以负端第一列为例, 对角分量为
\[
\dfrac{1}{z}\int_{-\infty}^{x}(\mathcal C_1)_{11}(y)\,\dd y
=\dfrac{\ii J_-(x)}{z},
\]
非对角分量为
\[
\dfrac{\ii}{z}\int_{-\infty}^{x}e^{2\ii\lambda(x-y)}\conj{q_y(y)}\,\dd y.
\]
由于 $q_x$ 在负端趋于零, 分部积分得
\[
\int_{-\infty}^{x}e^{2\ii\lambda(x-y)}\conj{q_y(y)}\,\dd y
=-\dfrac{\conj{q_x(x)}}{2\ii\lambda}
+\dfrac{1}{2\ii\lambda}\int_{-\infty}^{x}e^{2\ii\lambda(x-y)}\conj{q_{yy}(y)}\,\dd y.
\]
在 $\overline{D_+}$ 中, $y\le x$ 时指数模不超过 1, 且当 $z\to\infty$ 时 $|\lambda|\ge c|z|$. 因此该非对角分量绝对值至多为
\[
\dfrac{\|q_x\|_{L^\infty}+\|q_{xx}\|_{L^1}}{2|z\lambda|}=O(z^{-2}).
\]
正端和第二列只改变积分方向或指数符号, 具有相同估计. 于是
\[
n_\pm=I+\dfrac{\ii J_\pm\sg}{z}+O(z^{-2}),
\]
与 $\mathcal T$ 相乘即证明 \eqref{eq:mu-highenergy}. 此外, 从 \eqref{eq:n-volterra-system} 及 $n_\pm$ 非对角部分为 $O(z^{-2})$, 可得 $(n_\pm)_x=O(z^{-1})$, 因而 $(\mu_\pm)_x=O(z^{-1})$, 仍按解析列理解.

将展开代入 Wronskian 表达式, 有
\[
\det(\mu_{-,1},\mu_{+,2})
=1+\dfrac{\ii(J_--J_+)}{z}+O(z^{-2}),\qquad
J_--J_+=J_{\rm tot},\qquad \gamma=1+O(z^{-2}).
\]
这给出 $s_{22}$ 的展开. 同理,
\[
\det(\mu_{+,1},\mu_{-,2})
=1-\dfrac{\ii J_{\rm tot}}{z}+O(z^{-2}),
\]
从而得到 \eqref{eq:transmission-highenergy}. 在实轴上 $\theta$ 为实数, 由散射关系
\[
\begin{aligned}
S&=e^{-\ii\theta\sg}\mu_-^{-1}\mu_+e^{\ii\theta\sg},\\
\mu_-^{-1}\mu_+&=I+\dfrac{A_+-A_-}{z}+O(z^{-2})
=I-\dfrac{\ii J_{\rm tot}}{z}\sg+O(z^{-2}),
\end{aligned}
\]
得到 \eqref{eq:scattering-two-infinities} 的第一式.

最后, 令 $w=h(z)=-\dfrac{q_0^2}{z}$. 第 \ref{ch:symmetry} 节的换叶关系给出
\[
\mu_\pm(z)=-\dfrac{\ii q_0}{z}\mu_\pm(w)\sigma_1
=-\dfrac{\ii q_0}{z}\sigma_1+\dfrac{\ii}{q_0}A_\pm\sigma_1+O(z).
\]
计算 $\ii A_\pm\sigma_1=B\sigma_1-J_\pm\sg\sigma_1$ 即得 \eqref{eq:smallmu}. 因 $h$ 交换 $D_\pm$, 透射系数满足 $s_{22}(z)=s_{11}(w)$ 和 $s_{11}(z)=s_{22}(w)$; 将 $w^{-1}=-\dfrac{z}{q_0^2}$ 代入 \eqref{eq:transmission-highenergy}, 得到 \eqref{eq:transmission-smallenergy}. 在实轴上再用 $S(z)=\sigma_1S(w)\sigma_1$, 即得 \eqref{eq:scattering-two-infinities} 的第二式.
\end{proof}

两处无穷远控制 RHP 的归一化及光场重构. 分支点 $z=\pm\ii q_0$ 的连续延拓由命题 \ref{prop:jost-branch} 给出; 轮廓交点 $z=\pm q_0$ 是不同谱弧的交汇处, 共振原像 $z_a^\pm$ 则来自传播 Lax 矩阵的有理项.
